% Chapter 3 (two-cause rewrite, 2026-11). No result number in this chapter: results are in Ch. 6.
% Old text: chapters_old/chapter3.tex (git tag thesis-pre-two-causes).
\chapter{Methodology}
\label{ch:methodology}

This chapter presents how we study forgetting in continual \gls{kie} and how we address it.
In the following, we first formalise the problem and the metrics
(Section~\ref{sec:m-formulation}). Then, we describe how we analyse the two causes of
forgetting that this thesis separates: the drift of the representation
(Section~\ref{sec:m-drift}) and the conflict on the background label
(Section~\ref{sec:m-conflict}). After that, we detail the proposed method, CoLaR-SBR, and its
variant SBR-Recall (Sections~\ref{sec:m-overview} to~\ref{sec:m-recall}). Finally, we account
for its memory and compute cost (Section~\ref{sec:m-accounting}). All results are reported in
Chapter~\ref{ch:results}.

% =============================================================================
\section{Problem Formulation and Metrics}
\label{sec:m-formulation}

\paragraph{Continual key-information extraction.} A document is defined as
$x=(w_{1:L},\,b_{1:L},\,I)$, where $w_{1:L}$ are its \gls{ocr} tokens, $b_{1:L}$ their
normalised bounding boxes and $I$ the page image. A model $f_\theta$ labels every token with
a \gls{bio} tag, and the entities are scored by entity-level F1~\citep{tjong2003conll}. The
model consists of an encoder, which maps the document to one feature vector $h(x_i)\in
\mathbb{R}^{d}$ per token, and a linear head, which maps the feature to the tag logits
$z(x_i)$.

The documents arrive in $T$ sessions. Session $t$ brings a training set $\mathcal{D}_t$ that
is annotated with its own tag set $\mathcal{T}_t$ only. Specifically, a token whose true tag
is $y^{\ast}$ is labelled
\begin{equation}
  y=\begin{cases} y^{\ast} & y^{\ast}\in\mathcal{T}_t,\\ \mathrm{O} & \text{otherwise.}\end{cases}
  \label{eq:mask}
\end{equation}
In the \gls{bio} scheme, ``O'' is the tag of every token that belongs to no annotated entity.
From now on, we call it the \emph{background label}: we write ``O'' when we refer to the tag in
the data, and background label when we refer to its role in training. Under
Equation~\eqref{eq:mask}, the background label of session $t$ covers two kinds of tokens: the
tokens that are truly not entities, and the entities that session $t$ does not annotate. While
the model learns session $t$, it has no access to the training sets of the earlier sessions,
except for what a bounded memory keeps. After the last session, the model is tested on the
classes of all sessions, without knowing which session a document comes from.

\paragraph{Three settings.} We study three settings, which differ in what changes between
two sessions. In \gls{cil}, the tag set grows and the documents stay in one domain. This is
the setting where the background label changes its meaning from one session to the next;
hence, it is where our main claims are tested. In \gls{dil}, the tag set is fixed and the
documents change, e.g., from forms to receipts. Since every session annotates all classes,
Equation~\eqref{eq:mask} masks nothing; thus, forgetting can only come from the change of the
model itself, and \gls{dil} serves as the control without the conflict. In the \emph{Mixed}
setting, both change, which tests the two causes together. The six streams that instantiate
these settings on five public datasets are described in Chapter~\ref{ch:experimental-setup}.

\paragraph{Retention matrix.} To measure what the model keeps, we hold one test split per
session, masked like its training set. Let $R[i,j]$ be the F1 on the test split of session
$j$ after the model has been trained on session $i$. From this matrix, we report the
\gls{aa} and the \gls{af}:
\begin{equation}
  \mathrm{AA}=\frac{1}{T}\sum_{j=1}^{T}R[T,j],\qquad
  \mathrm{AF}=\frac{1}{T-1}\sum_{j=1}^{T-1}\big(R[j,j]-R[T,j]\big).
  \label{eq:aa-af}
\end{equation}
\gls{aa} is the mean F1 of the final model over the sessions. \gls{af} is the F1 of a session
right after it was learned minus its final F1, averaged over the sessions before the last
one. It is noted that the diagonal $R[j,j]$ matters on its own: it tells whether a session
was learned at all, which a single final number cannot show.

\paragraph{Session-restricted scoring.} The test labels of session $j$ show only its own
classes. Hence, raw scoring would charge a model for every correct entity of another
session: such an entity is labelled ``O'' in the split of session $j$, so predicting it
counts as a false positive. To avoid this, we score each split \emph{session-restricted}: a
predicted tag outside $\mathcal{T}_j$ is mapped to ``O'' before the spans are matched. As a
result, a predicted entity can be a false positive only in its own session, and a missed
entity is counted only in its own session.

\paragraph{Full-label F1.} Session-restricted scoring has two limits. The test split of a
session keeps only the documents that contain one of its entities, and \gls{aa} weights all
sessions equally, although they hold different numbers of entities. Hence, on the
class-incremental streams we also report the \emph{full-label F1}: the entity-level F1 of
the final model on the complete test set with all classes and no label masked, as in
continual named entity recognition~\citep{ma2023learning}. The full-label F1 is our headline
number on the class-incremental streams, and \gls{aa} and \gls{af} describe retention on
every stream. Table~\ref{tab:notation} summarises the notation of this chapter.

\begin{table}[!htbp]
  \centering
  \caption[Notation of the methodology]{\textbf{Notation.} Symbols used in this chapter.}
  \label{tab:notation}
  \begin{tabular}{@{}ll@{}}
    \toprule
    Symbol & Meaning \\
    \midrule
    $T$, $t$ & number of sessions; index of the current session \\
    $\mathcal{D}_t$, $\mathcal{T}_t$ & training set and tag set of session $t$ \\
    $h(x)\in\mathbb{R}^{d}$, $z(x)$ & feature and tag logits of a token $x$ \\
    $R[i,j]$ & F1 on the test split of session $j$ after training session $i$ \\
    $k$, $r$ & layer where the memory is injected; rank of its compression \\
    $H_j\in\mathbb{R}^{L\times d}$ & layer-$k$ activations of the stored document $j$ \\
    $b$, $\mathcal{T}_b$ & a readout block and the tags it holds \\
    $o_b(x)$, $p_b(\cdot\mid x)$ & background logit and distribution of block $b$ \\
    $m_{sb}$ & whether a document of session $s$ supervises block $b$ \\
    $e_b(x)$ & claim of block $b$ on a token, $1-p_b(\mathrm{O}_b\mid x)$ \\
    $\theta_0$, $\Delta_b$ & frozen encoder state and the low-rank update of block $b$ \\
    \bottomrule
  \end{tabular}
\end{table}

% =============================================================================
\section{Analysing Cause 1: Representation Drift}
\label{sec:m-drift}

The first cause of forgetting is the drift of the model. When the encoder is trained on a
new session, the features that it computes for the documents of the old sessions change, and
the head that was fitted on the old features no longer reads them correctly. This cause does
not depend on the labels; it is present even when the label space is fixed. In this section,
we describe how we measure the drift and how we test what removes it.

\paragraph{Representational drift.} To measure how much the representation of a layer
changes, we use linear \gls{cka}~\citep{kornblith2019similarity}. Given two centred
activation matrices $X\in\mathbb{R}^{n\times p}$ and $Y\in\mathbb{R}^{n\times q}$ of the same
$n$ inputs,
\begin{equation}
  \mathrm{CKA}(X,Y)=\frac{\lVert Y^{\top}X\rVert_F^{2}}
  {\lVert X^{\top}X\rVert_F\,\lVert Y^{\top}Y\rVert_F}.
  \label{eq:cka}
\end{equation}
We compute it between the activations of the model before and after a session on a fixed
probe set; a value of $1$ means that the layer did not drift. Two choices matter here.
First, token classification reads one feature per token; hence, we compute \gls{cka} over
individual tokens and not over the mean vector of a document, which could report no drift
while the features of the tokens have been reorganised. Second, the probe set holds more
tokens than the width of the layer, so that the estimate is stable.

\paragraph{Importance and displacement.} \gls{cka} tells where the features change. To tell
where the \emph{parameters} that mattered for an old session are overwritten, we use two
quantities per parameter group $g$. The empirical Fisher importance on the data
$\mathcal{D}$ of a session,
\begin{equation}
  F_g=\frac{1}{|\mathcal{D}|}\sum_{(x,y)\in\mathcal{D}}\;\sum_{\phi\in g}
  \Big(\frac{\partial\log p_\theta(y\mid x)}{\partial\phi}\Big)^{2},
  \label{eq:fisher}
\end{equation}
measures how sensitive the loss of that session is to the group. It is noted that importance
alone is not a sign of forgetting: the head carries the largest gradients of any token
classifier, whether or not its parameters are moved later. Hence, we weight the change of the
parameters between two checkpoints by the importance of the old session,
\begin{equation}
  D_g=\sum_{\phi\in g}F^{(t-1)}_{\phi}\,\big(\theta^{(t)}_{\phi}-\theta^{(t-1)}_{\phi}\big)^{2},
  \label{eq:displacement}
\end{equation}
which is large exactly where parameters that mattered before session $t$ have been changed
by session $t$. We report the three quantities by depth, in five groups: the input
embeddings, the early, middle and late layers, and the head.

\paragraph{Protocol.} We fine-tune the model naively on a stream, i.e., one session after the
other with no continual-learning method, and keep a checkpoint at every session boundary.
At each boundary, we compute \gls{cka}, the Fisher importance and the displacement between
the two checkpoints, together with one row of the retention matrix
(Algorithm~\ref{alg:drift}). To examine whether the picture depends on the encoder, we repeat
the analysis on four encoders: LayoutLMv3~\citep{huang2022layoutlmv3}, which reads text,
layout and the page image; LiLT~\citep{wang2022lilt} and BROS~\citep{hong2022bros}, which
read text and layout; and BERT~\citep{devlin2019bert}, which reads text only. On LayoutLMv3,
we additionally repeat it with some of the inputs masked (text only, text and layout, image
and layout), so that the architecture is fixed and only the inputs differ.

\begin{algorithm}[!htbp]
  \caption{Drift analysis at the session boundaries}
  \label{alg:drift}
  \begin{algorithmic}[1]
    \Require sessions $\mathcal{D}_{1:T}$, probe set $\mathcal{P}$, layers $\mathcal{G}$, parameter groups $\mathcal{G}_F$
    \State initialise $\theta^{(0)}$ from the pre-trained encoder
    \For{$t=1$ \textbf{to} $T$}
      \State grow the head by the tags $\mathcal{T}_t$
      \State $\theta^{(t)}\gets\textsc{FineTune}(\theta^{(t-1)},\mathcal{D}_t)$ \Comment{naive, no memory}
      \ForAll{$\ell\in\mathcal{G}$}
        \State $\mathrm{CKA}_{t,\ell}\gets\textsc{LinearCKA}\big(\textsc{Act}_\ell(\theta^{(t-1)},\mathcal{P}),\,\textsc{Act}_\ell(\theta^{(t)},\mathcal{P})\big)$
      \EndFor
      \ForAll{$g\in\mathcal{G}_F$}
        \State $F_{t,g}\gets\textsc{Fisher}(\theta^{(t)},g,\mathcal{D}_t)$;\quad
               $D_{t,g}\gets\sum_{\phi\in g}F^{(t-1)}_{\phi}(\theta^{(t)}_{\phi}-\theta^{(t-1)}_{\phi})^{2}$
      \EndFor
      \State $R[t,j]\gets\textsc{F1}(\theta^{(t)},\text{test split of session } j)$ for all $j\le t$
    \EndFor
    \State \Return $\mathrm{CKA}$, $F$, $D$, $R$
  \end{algorithmic}
\end{algorithm}

\paragraph{What removes the drift.} The measurement tells where the drift sits, not what
removes it. To examine the latter, we compare the main families of continual learning on the
domain-incremental stream, where drift is the only cause: regularisation of the
parameters~\citep{kirkpatrick2017ewc}, distillation from the previous
model~\citep{li2017lwf} and replay of old documents~\citep{rolnick2019er,buzzega2020der},
against naive fine-tuning and joint training. If drift is what a method must counter, then a
method that shows old documents to the model again should come close to joint training on
this stream, whereas a method that only constrains the parameters may not. This comparison is
repeated on the four encoders (Section~\ref{sec:c6-s1}), and its outcome decides the starting
point of our method.

% =============================================================================
\section{Analysing Cause 2: The Background-Label Conflict}
\label{sec:m-conflict}

The second cause appears only when the tag set grows. In this section, we first show where
the conflict comes from, and then state the series of hypotheses with which we test it.

\subsection{Where the Conflict Comes From}
\label{sec:m-conflict-origin}

Consider the usual readout, one softmax over the background label and all classes seen so
far,
\begin{equation}
  p(c\mid x)=\frac{\exp z_c(x)}{\exp z_{\mathrm{O}}(x)+\sum_{c'\in\mathcal{T}_{1:t}}\exp z_{c'}(x)},
  \qquad c\in\{\mathrm{O}\}\cup\mathcal{T}_{1:t},
  \label{eq:shared}
\end{equation}
where $\mathcal{T}_{1:t}=\bigcup_{s\le t}\mathcal{T}_s$. We refer to it as the \emph{shared
softmax}. For a token that is labelled ``O'', the gradient of the cross-entropy with respect
to the logits is
\begin{equation}
  \frac{\partial\mathcal{L}}{\partial z_{\mathrm{O}}}=p(\mathrm{O}\mid x)-1\le 0,\qquad
  \frac{\partial\mathcal{L}}{\partial z_c}=p(c\mid x)\ge 0\quad\text{for every } c\in\mathcal{T}_{1:t}.
  \label{eq:grad}
\end{equation}
Thus, one background token raises the background logit and lowers the logit of \emph{every}
class in the head, including the classes that its session did not annotate. Under
Equation~\eqref{eq:mask}, this happens on two sides (Figure~\ref{fig:conflict}).

\begin{itemize}
  \item \textbf{Current side.} A document of the current session $t$ labels every entity of
        an earlier session as ``O''. Hence, each training step lowers the logits of the old
        classes on tokens that truly belong to them.
  \item \textbf{Stored side.} A document that was stored at session $s<t$ labels every entity
        of a later session as ``O''. Hence, each replay step lowers the logits of the new
        classes on tokens that truly belong to them. The stored labels are \emph{stale}: they
        were complete when the document was stored and became wrong when the head grew.
\end{itemize}

\noindent As a result, the same token receives contradicting labels over time, and the
single background logit $z_{\mathrm{O}}$ is trained by documents that disagree on what the
background is. It is noted that a memory does not remove this conflict. On the contrary,
replay adds the stored side to it; this is why we treat the conflict as a cause of its own
and not as a harder form of drift. The shift of the background label is known from continual
semantic segmentation~\citep{cermelli2020mib} and continual named entity
recognition~\citep{zhang2023cpfd,ma2023learning}, and its two sides were named in continual
object detection~\citep{liu2023abr,mo2024bpf}. However, it is unclear which side matters in
document streams, and whether repairing the labels is enough.

\subsection{A Series of Hypotheses}
\label{sec:m-hypotheses}

We examine the failure of replay under label growth with three hypotheses. Each one makes a
prediction, and each is tested by changing one intervention against the run it is compared
with (Table~\ref{tab:hypotheses}). For each hypothesis, we report whether its prediction
holds, fails or holds in part.

\paragraph{H1: only harder forgetting.} We first hypothesize that label growth only makes
forgetting harder, i.e., that the failure is still caused by drift. If so, a memory that
works under domain shift should also work under label growth, and a larger memory should
close the remaining gap. To examine it, we take the same memory with the shared softmax from
the domain-incremental stream to the class-incremental one, we vary its size, and we replace
the stored activations by raw documents. We also read the diagonal of the retention matrix:
under drift alone, a session is learned well and lost later, whereas under the conflict the
session that is being trained may not be learned in the first place.

\paragraph{H2: the labels are the cause.} We next hypothesize that the labels are the cause,
on the two sides above. If so, repairing the labels should close the gap with the shared
softmax unchanged. On the stored side, we repair the labels with the ground-truth labels,
which a continual learner does not have, in two ways: every stored document receives the
true labels of all classes in the grown head before each session; or the stale tokens are
masked in the replay loss, which removes the false background label without adding a label.
On the current side, we use two repairs that need no additional annotation. The first one
follows the unbiased cross-entropy of MiB~\citep{cermelli2020mib}: a background token is
allowed to be the background \emph{or} any class that its document could not annotate,
\begin{equation}
  \mathcal{L}_{\mathrm{MiB}}(x,\mathrm{O},s)=-\log\Big(p(\mathrm{O}\mid x)
  +\sum_{c\,\in\,\mathcal{T}_{1:t}\setminus\mathcal{T}_s}p(c\mid x)\Big),
  \label{eq:mib}
\end{equation}
where $s$ is the session of the document, so that the sum runs over the earlier classes for
a current document and over the later classes for a stored one. The second one slices the
softmax: the loss of a stored document is normalised over the classes that were known when it
was stored, and, in a further control inspired by SS-IL~\citep{ahn2021ssil}, the loss of the
current background tokens is normalised without the old classes.

\paragraph{H3: the shared normalisation.} Finally, we hypothesize that the conflict lives in
the shared normalisation itself: one background logit that is trained by documents which
disagree on the background. If so, giving each session its own background logit, normalised
with the tags of that session only, should close most of the gap without any additional
annotation. This is the diagnosis of SS-IL for a shared softmax~\citep{ahn2021ssil} with a
background label added, and the readout of ICE-O~\citep{liu2023ice} without its frozen
encoder. We test it with and without the memory, in order to separate what the readout
contributes from what replay contributes.

\begin{table}[!htbp]
  \centering
  \caption[The series of hypotheses and their tests]{\textbf{The series of hypotheses.}
  Each hypothesis is tested by one change against the memory with a shared softmax.
  $^{r}$Uses the ground-truth labels of the stored documents.}
  \label{tab:hypotheses}
  \begin{tabularx}{\textwidth}{@{}l>{\raggedright\arraybackslash}p{4.3cm}L@{}}
    \toprule
    Hypothesis & Intervention & Prediction if the hypothesis holds \\
    \midrule
    H1 only harder forgetting & same memory as on \gls{dil}; larger memory; raw documents &
      the gap to joint training closes as on \gls{dil}; every session is learned when it is trained \\
    \addlinespace
    H2 the labels (stored side) & true labels before each session$^{r}$; stale tokens masked$^{r}$ &
      the shared softmax reaches joint training on every stream \\
    H2 the labels (current side) & MiB-style loss (Equation~\ref{eq:mib}); softmax slices &
      the gap closes without ground-truth labels \\
    \addlinespace
    H3 the shared normalisation & one background logit per session, with and without the memory &
      most of the gap closes on every stream, with no additional annotation \\
    \bottomrule
  \end{tabularx}
\end{table}

\noindent The three hypotheses are tested on the three class-incremental streams in
Section~\ref{sec:c6-s2}. The outcome of H3 motivates the readout of our method.

% =============================================================================
\section{Overview of the Proposed Method}
\label{sec:m-overview}

In this work, we aim to remove the two causes: the drift of the encoder and the conflict on
the background label. The workflow of the proposed method, CoLaR-SBR, is shown in
Figure~\ref{fig:method}. Given the documents of a session, the model is trained on them as
usual. After the session, we store a few of its documents as compressed activations of an
intermediate layer (step~1). Then, while the model learns a new session, we replay the
stored documents to remind the model of the past sessions (step~2); this addresses the
drift. Finally, every session reads its tokens with its own background logit (step~3), so
that the background label of one session is never counted against the classes of another;
this addresses the conflict. A variant, SBR-Recall, keeps one low-rank update of the encoder
per session in place of the memory.

\begin{figure}[!htbp]
  \centering
  \papertikz{0.95}{figures/paper/tikz_method}
  \caption[Overview of CoLaR-SBR]{\textbf{Overview of CoLaR-SBR.} Documents are (1) stored
  and (2) replayed as layer-$k$ activations; (3) each session reads its tokens with its own
  background logit (b), not with a shared one (a). SBR-Recall keeps one low-rank update per
  session and no memory.}
  \label{fig:method}
\end{figure}

\paragraph{The principle.} We argue that a token labelled as background in session $s$ means
``none of the classes annotated in session $s$''. It is evidence against those classes and
silent about every other class, as in training on partially annotated
data~\citep{greenberg2018marginal}. Hence, we \emph{supervise each class only with negatives
from documents that annotated it, and normalise together only the classes that were annotated
together}. The repairs of H2 follow the same principle on the side of the labels or of the
loss; our readout follows it in the normalisation, where no label needs to be repaired.

% =============================================================================
\section{CoLaR: Compressed Latent Replay}
\label{sec:m-colar}

The readout of Section~\ref{sec:m-sbr} needs stored documents that keep the label of each
token at its position, together with the session that wrote the labels. Hence, our memory
keeps \emph{whole documents}, not statistics of classes. Inspired by latent
replay~\citep{hayes2020remind,pellegrini2020latent}, CoLaR (Compressed Latent Replay) stores a
document as the activations of an intermediate layer and not as its raw page. It works in
three steps.

\paragraph{Freezing the lower layers.} First, after the first session, we freeze the input
embeddings and the layers $1..k$ of the encoder. A stored activation stays valid only while
the layers below it do not change; with these layers frozen, an activation that was stored at
any session is exactly what the current model would compute for the same document.

\paragraph{Storing.} Then, at the end of each session, we draw $N$ documents at random from
its training set. For each document $j$, we take its layer-$k$ activations
$H_j\in\mathbb{R}^{L\times d}$, where $L$ is the number of positions of the document, and
compress them with a truncated \gls{svd} of rank $r$,
\begin{equation}
  H_j\approx U_j\Sigma_jV_j^{\top},\qquad U_j\in\mathbb{R}^{L\times r},\;
  \Sigma_j\in\mathbb{R}^{r\times r},\;V_j\in\mathbb{R}^{d\times r}.
  \label{eq:svd}
\end{equation}
We store the factors $U_j\Sigma_j$ and $V_j$ in half precision, with the boxes, the attention
mask, the labels and the session of the document. It is noted that the compression is
applied to one document at a time. The features of all documents together are close to full
rank, whereas the activations of a single document are not, since its tokens share one
layout and one vocabulary; thus, a per-document factorisation keeps the document faithfully
at a rank far below $d$. No pixel and no token id is kept.

\paragraph{Replaying.} Finally, at every training step, we draw stored documents, rebuild
their activations $H_j$ from the factors and inject them at layer $k{+}1$, where the boxes
and the mask of the document drive the layout biases and the attention mask as for a current
document. The upper layers and the head are then trained on the current batch and on the
replayed documents together. We use $k{=}4$, $r{=}128$ and $N{=}50$ documents per session;
the impact of each setting is studied in Section~\ref{sec:c6-s5}.

\paragraph{Why activations and not raw documents.} A memory of raw documents would serve the
readout as well, and we compare both in Section~\ref{sec:c6-s4}. We choose the latent memory
for two reasons. First, it holds neither pixels nor token ids. Second, a replayed document
passes only through the layers above $k$; hence, replay trains fewer parameters than raw
replay does. We do not claim that it is smaller than a raw document in every encoding
(Section~\ref{sec:m-accounting}).

% =============================================================================
\section{SBR: Session-Blocked Readout}
\label{sec:m-sbr}

SBR (Session-Blocked Readout) is our instantiation of a readout with one background per
session~\citep{liu2023ice,kim2022more}. In the following, we describe its blocks, its
training rule and its inference rule.

\paragraph{Blocks.} Every session that brings new tags opens a \emph{block} $b$ that holds
those tags $\mathcal{T}_b$. Whereas the shared softmax of Equation~\eqref{eq:shared}
normalises over the background label and all classes, SBR gives every block its own
background logit, a linear map of the token feature,
\begin{equation}
  o_b(x)=w_b^{\top}h(x)+c_b,
  \label{eq:olog}
\end{equation}
and normalises each block separately:
\begin{equation}
  p_b(\cdot\mid x)=\operatorname{softmax}\big([\,o_b(x),\,z_{\mathcal{T}_b}(x)\,]\big),
  \label{eq:block}
\end{equation}
where $z_{\mathcal{T}_b}(x)$ are the entries of the tag logits for the tags in
$\mathcal{T}_b$. We write $\mathrm{O}_b$ for the background of block $b$. A session that only
re-annotates known tags opens no block; hence, SBR has a single block on a fixed label
space. The tag logits come from the usual linear head, so SBR adds only one vector $w_b$ and
one bias $c_b$ per session.

\paragraph{Training.} For a token $x$ with label $y$ in a document of session $s$, the loss
sums the cross-entropy of the blocks that the document can supervise:
\begin{equation}
  \mathcal{L}(x,y,s)=\sum_{b}m_{sb}\,\mathrm{CE}\big(p_b(\cdot\mid x),\,y_b\big),\qquad
  y_b=\begin{cases} y & y\in\mathcal{T}_b,\\ \mathrm{O}_b & \text{otherwise.}\end{cases}
  \label{eq:sbr}
\end{equation}
The mask $m_{sb}$ decides whether the labels of the document say anything about block $b$.
It follows three cases.
\begin{enumerate}
  \item \emph{Full.} $m_{sb}=1$ when session $s$ annotates every tag of $\mathcal{T}_b$. Its
        background labels are then true negatives for the classes of block $b$, and its
        entity labels are positives.
  \item \emph{All-background.} $m_{sb}=1$ with $y_b=\mathrm{O}_b$ on every token when block
        $b$ belongs to a different annotation schema than session $s$, e.g., a block of form
        fields and a receipt. No entity of block $b$ can occur in such a document; hence,
        every token is a negative for $b$.
  \item \emph{Skipped.} $m_{sb}=0$ otherwise, because the background labels of the document
        say nothing about block $b$: an ``O'' may hide one of its entities.
\end{enumerate}
A stored document follows the rule of the session that stored it, which is why the memory
keeps the session of each document. Table~\ref{tab:supervision} gives examples. On a
class-incremental stream within one schema, a document trains only the block of its own
session, so neither a current nor a stored background label can count against the classes of
another session; this removes both sides of the conflict of
Section~\ref{sec:m-conflict-origin}. Across schemas, the all-background case keeps a block
from claiming the tokens of another dataset. It is noted that the shared softmax is the
special case of a single block that every session supervises.

\begin{table}[!htbp]
  \centering
  \caption[The supervision rule of SBR on examples]{\textbf{The supervision rule of SBR.}
  Which blocks a document trains, and how. ``--'': the block does not exist yet or is
  skipped.}
  \label{tab:supervision}
  \begin{tabularx}{\textwidth}{@{}>{\raggedright\arraybackslash}p{4.4cm}>{\raggedright\arraybackslash}p{2.6cm}>{\raggedright\arraybackslash}p{2.9cm}L@{}}
    \toprule
    Document & Own block & Other blocks, same schema & Blocks of another schema \\
    \midrule
    \multicolumn{4}{@{}l}{\emph{Class-incremental stream (one schema)}} \\
    Current document of session 3 & block 3: full & blocks 1, 2: skipped & -- \\
    Stored document of session 1, replayed in session 3 & block 1: full & blocks 2, 3: skipped & -- \\
    \addlinespace
    \multicolumn{4}{@{}l}{\emph{Mixed stream (three schemas)}} \\
    Receipt of a new dataset & its block: full & -- & form blocks: all-background \\
    Form of the last session, all form labels & both form blocks: full & -- & receipt blocks: all-background \\
    \addlinespace
    \multicolumn{4}{@{}l}{\emph{Domain-incremental stream (fixed label space)}} \\
    Any document & the single block: full & -- & -- \\
    \bottomrule
  \end{tabularx}
\end{table}

\paragraph{Inference.} At test time, each block reports its \emph{claim} on a token, the
probability that the token is one of its entities,
\begin{equation}
  e_b(x)=1-p_b(\mathrm{O}_b\mid x).
  \label{eq:claim}
\end{equation}
The block with the largest claim labels the token with its best tag if the claim passes a
threshold $\tau$; otherwise, the token is ``O'':
\begin{equation}
  \hat y(x)=\begin{cases}\arg\max_{c\in\mathcal{T}_{b^{\star}}}p_{b^{\star}}(c\mid x) & e_{b^{\star}}(x)>\tau,\\[2pt]
  \mathrm{O} & \text{otherwise,}\end{cases}\qquad b^{\star}=\arg\max_b e_b(x).
  \label{eq:infer}
\end{equation}
Hence, no session identity is needed at test time: the blocks decide among themselves which
one owns a token. We set $\tau=0.5$ and do not tune it. Since each block is a softmax of its
tags against its own background, $e_b>0.5$ means that the block finds an entity more likely
than its background, which is the decision the block was trained to make.
Algorithm~\ref{alg:colarsbr} summarises one session of CoLaR-SBR.

\begin{algorithm}[!htbp]
  \caption{One session of CoLaR-SBR}
  \label{alg:colarsbr}
  \begin{algorithmic}[1]
    \Require session $t$ with training set $\mathcal{D}_t$ and tags $\mathcal{T}_t$; memory $\mathcal{M}$; layer $k$, rank $r$, $N$ documents
    \If{$\mathcal{T}_t$ holds new tags} open block $t$: grow the head by $\mathcal{T}_t$, add $(w_t,c_t)$ \EndIf
    \If{$t=2$} freeze the embeddings and the layers $1..k$ \EndIf
    \For{each training step}
      \State draw a batch $B\subset\mathcal{D}_t$ and stored documents $B_{\mathcal{M}}\subset\mathcal{M}$
      \State $\mathcal{L}_{\text{cur}}\gets$ mean of Equation~\eqref{eq:sbr} over the tokens of $B$, with $s=t$
      \State rebuild $H_j=U_j\Sigma_jV_j^{\top}$ for $j\in B_{\mathcal{M}}$ and run the layers $k{+}1..12$ on them
      \State $\mathcal{L}_{\text{rep}}\gets$ mean of Equation~\eqref{eq:sbr} over their tokens, with $s$ the session that stored $j$
      \State update the trainable parameters with $\mathcal{L}_{\text{cur}}+\mathcal{L}_{\text{rep}}$
    \EndFor
    \State draw $N$ documents of $\mathcal{D}_t$; store their rank-$r$ factors of layer $k$ with boxes, mask, labels and $t$ in $\mathcal{M}$
  \end{algorithmic}
\end{algorithm}

\paragraph{Relation to prior readouts.} Separated softmaxes were proposed for image
classification, where SS-IL~\citep{ahn2021ssil} normalises old and new classes apart and
ER-ACE~\citep{caccia2022erace} restricts the loss of the current batch to its classes. In
these settings there is no background label. ICE-O~\citep{liu2023ice} adds a constant
background logit to each step of a class-incremental tagger with a frozen encoder, and
MORE~\citep{kim2022more} trains one head per task with an out-of-distribution output.
Unlike them, SBR learns the background logit of each block from the token feature, trains
the encoder with a memory, and decides from the annotation of a document which blocks it may
supervise. The contribution of each of these choices is measured in
Section~\ref{sec:c6-s4}.

% =============================================================================
\section{SBR-Recall: Per-Session Deltas}
\label{sec:m-recall}

SBR-Recall is a variant of our method that keeps no memory. Without the conflict, the memory
is needed only against drift, and the task-identity view of class-incremental
learning~\citep{kim2022theoretical} states when it can be removed: if each block is read
under the encoder state that learned it, nothing drifts for that block, provided that the
model can tell which state a token needs. SBR-Recall realises this in three steps.

\paragraph{Per-session deltas.} First, the model is trained on the first session as usual,
and the encoder is then frozen; we denote this state by $\theta_0$. Then, for each later
session $b$, we train a low-rank update $\Delta_b$~\citep{hu2022lora} of the layers above
$k$, together with the head and the background logits, and store it ($\Delta_1=0$). Every
linear layer of the attention and of the feed-forward part carries one such update per
session, initialised so that a new update starts as the identity.

\paragraph{Reading each block under its own state.} Finally, every term of the loss and
every claim of block $b$ is computed under the state that learned the block:
\begin{equation}
  p_b(\cdot\mid x)=\operatorname{softmax}\big([\,o_b,\,z_{\mathcal{T}_b}\,]\big)
  \big(x;\,\theta_0+\Delta_b\big).
  \label{eq:recall}
\end{equation}
During training, the blocks that a document supervises are grouped by their update, and one
forward pass is run per update, so that each term is computed under the encoder state of its
own block. At test time, the claims $e_b$ of Equation~\eqref{eq:claim} are computed in the
same way and select the block, and with it the state, by Equation~\eqref{eq:infer}. Since the
state $\theta_0+\Delta_b$ never changes after session $b$, the predictions of block $b$ on a
given token do not change either; a later session can only change which block wins.

\paragraph{When the memory can be removed.} This selection relies on the background logit of
each block being trained on true negatives. Under label growth within one schema, this
holds: the documents of a session contain the entities of the other sessions labelled ``O'',
so each block learns to reject them. In contrast, across schemas or on a fixed label space, a
block never sees the tokens of the other sessions during its own session; hence, we expect
SBR-Recall to need the stored documents there, in which case the memory of
Section~\ref{sec:m-colar} is kept beside the updates. Both cases are examined in
Sections~\ref{sec:c6-s3} and~\ref{sec:c6-s4}.

% =============================================================================
\section{Memory and Compute Accounting}
\label{sec:m-accounting}

In this section, we state how the cost of the method is counted, so that the comparisons of
Section~\ref{sec:c6-s7} can be checked.

\paragraph{Memory of CoLaR.} A stored document holds the two factors in half precision and
its boxes, mask and labels. With $L$ positions and width $d$, the factors take $2r(L+d)$
bytes. For instance, LayoutLMv3 has $L=709$ positions ($512$ text tokens and $197$ image
patches) and $d=768$; with $r=128$, the factors take $2\cdot128\cdot(709+768)=378{,}112$
bytes, and the boxes, mask and labels add $24{,}576$ bytes, i.e., about $0.4$\,MB per
document. A stream of $T$ sessions with $N$ documents each therefore stores $TN$ documents,
plus one byte per document for its session under SBR. For comparison, the uncompressed
activations of a document take $2Ld$ bytes, and a raw document as our code holds it takes
$3\cdot224^{2}$ pixels at four bytes each plus its tokens; we also report the size of a raw
document with 8-bit pixels, an encoding that we did not run.

\paragraph{Memory of SBR-Recall.} A rank-$\rho$ update of a linear layer with $n_{\text{in}}$
inputs and $n_{\text{out}}$ outputs has $\rho\,(n_{\text{in}}+n_{\text{out}})$ parameters. One
encoder layer carries updates on its query, key, value and attention-output maps
($d\rightarrow d$) and on its two feed-forward maps ($d\rightarrow 4d$ and $4d\rightarrow
d$); hence, an update over the $12-k$ upper layers has
\begin{equation}
  (12-k)\,\big[\,4\cdot 2\rho d+2\cdot 5\rho d\,\big]=18\,\rho\,d\,(12-k)
  \label{eq:delta-params}
\end{equation}
parameters. The cost grows with the number of sessions and not with the number of their
documents; thus, the updates are cheaper than a memory only when the documents of a session
take more than its update.

\paragraph{Training cost.} After the first session, CoLaR-SBR trains the layers above $k$,
the head and the background logits; the lower layers are frozen. A method that replays raw
documents trains the whole network. SBR-Recall trains one update, the head and the
background logits at a time.

\paragraph{Inference cost.} CoLaR-SBR needs one forward pass; the blocks share the encoder
and differ only in the last linear map. SBR-Recall needs one pass per session, since each
block is read under its own state. As the lower $k$ layers are shared by all states, an
implementation that caches them repeats only the $12-k$ upper layers; with $T$ sessions, the
cost is $1+(T-1)(12-k)/12$ times that of one pass.

% =============================================================================
\section{Summary}
\label{sec:m-summary}

This chapter separated forgetting in continual \gls{kie} into two causes and assigned one
instrument and one component to each (Table~\ref{tab:mapping}). The drift of the
representation is measured by \gls{cka}, importance and displacement at the session
boundaries, and is addressed by replaying whole documents from a compressed latent memory
(CoLaR), or removed by reading each block under its own frozen state (SBR-Recall). The
conflict on the background label is tested by a series of three hypotheses, and is addressed
at the readout, where each session receives its own background logit (SBR). The next two
chapters describe the implementation and the experimental setup, and
Chapter~\ref{ch:results} reports the outcome of each analysis and the comparison with the
baselines.

\begin{table}[!htbp]
  \centering
  \caption[Causes, tests and components]{\textbf{From the causes to the method.} Each cause
  has its test and the component that answers it.}
  \label{tab:mapping}
  \begin{tabularx}{\textwidth}{@{}lLL@{}}
    \toprule
    Cause & How it is tested & Component that answers it \\
    \midrule
    Drift of the representation & \gls{cka}, importance and displacement by depth; replay
      against regularisation on \gls{dil} & CoLaR: whole documents replayed as compressed
      activations; SBR-Recall: one frozen state per session \\
    \addlinespace
    Conflict on the background label & H1 to H3 on three class-incremental streams &
      SBR: one background logit per session, supervised only by the documents that
      annotated its classes \\
    \bottomrule
  \end{tabularx}
\end{table}
